\subsection{基线校验与占用热点识别}

进行关键资源敏感度分析之前，我们首先固定当前实例的\textbf{真实调度结果}。基线工况包含当前仓库布局、12 个静态任务、4 台 AMR，以及已注入的动态事件。系统依次完成\textbf{静态排程、冲突修复和动态重排}；同时将轨迹点投影到网格单元，统计\textbf{轨迹占用热点}。关键结果汇总见表~\ref{tab:baseline-summary} 和表~\ref{tab:baseline-hotspots}。

\begin{center}
\begin{minipage}[t]{0.53\linewidth}
\centering
\small
\captionof{table}{基线工况关键结果}
\label{tab:baseline-summary}
\renewcommand{\arraystretch}{1.76}
\begin{tabular}{@{}lll@{}}
\toprule
模块 & 任务 / AMR & $\Cmax$ \\
\midrule
静态排程 & 12 / 4 & 40.840111 \\
冲突修复 & 12 / 4 & 46.596970 \\
动态重排 & 13 / 4 & 61.918082 \\
\bottomrule
\end{tabular}
\renewcommand{\arraystretch}{1}
\end{minipage}
\hfill
\begin{minipage}[t]{0.39\linewidth}
\centering
\small
\captionof{table}{基线轨迹占用热点}
\label{tab:baseline-hotspots}
\begin{tabular}{@{}ccc@{}}
\toprule
热点 & 网格坐标 & 占用次数 \\
\midrule
1 & $(6,4)$ & 166 \\
2 & $(5,4)$ & 152 \\
3 & $(9,1)$ & 95 \\
4 & $(9,2)$ & 84 \\
5 & $(8,2)$ & 77 \\
\bottomrule
\end{tabular}
\end{minipage}
\end{center}

由表~\ref{tab:baseline-summary} 可见，\textbf{冲突修复使 $\Cmax$ 从 40.840111 上升到 46.596970，增幅约 14.1\%}；\textbf{动态重排后 $\Cmax$ 进一步上升到 61.918082，较冲突修复结果增加约 32.9\%}。与此同时，总延期、剩余冲突和硬违规均为 0，说明\textbf{基线排程具有可行性}；变更任务 4 个、新增任务 1 个，则说明\textbf{动态事件已经带来可观的计划扰动}。表~\ref{tab:baseline-hotspots} 显示，$(6,4)$ 与 $(5,4)$ 两个相邻网格占用最高，分别达到 166 次和 152 次；$(9,1)$、$(9,2)$ 和 $(8,2)$ 也形成\textbf{连续的高占用区域}。

为解释完工时间上升背后的空间原因，我将轨迹点投影到\textbf{网格单元} $\Omega_{ab}$，并统计\textbf{占用次数与归一化占用密度}：
\begin{equation}
H_{ab}=\sum_{k=1}^{N}\indicator\{(x_k,y_k)\in\Omega_{ab}\},
\qquad
\rho_{ab}=\frac{H_{ab}}{\sum_{p,q}H_{pq}}.
\end{equation}
对应空间分布见图~\ref{fig:baseline-dashboard}。图~\ref{fig:baseline-dashboard} 集中可视化展现了前文的\textbf{运筹策略与资源占有情况}.

\begin{center}
\includegraphics[width=0.94\linewidth]{baseline_dashboard}
\captionof{figure}{基线工况下的重排甘特图与轨迹占用热图}
\label{fig:baseline-dashboard}
\end{center}

综合表~\ref{tab:baseline-summary}、表~\ref{tab:baseline-hotspots} 和图~\ref{fig:baseline-dashboard}，\textbf{当前系统已经满足可行性校验，同时表现出局部拥堵与动态扰动}。敏感度分析因此围绕\textbf{AMR 资源冗余、任务负荷边界和瓶颈通道容量}展开。
